\makeatletter   % Para poder usar @ en los nombres de comandos

% Comando para las listas
\newcommand{\lista}[1]{
    \begin{minipage}[t]{10cm}
    \begin{itemize}[leftmargin=*, itemsep=0pt]
        #1
    \end{itemize}
    \break
    \end{minipage}
}

% Comando de tabla genérica
\define@key{tablarara}{titulo}{\def\tablarara@titulo{#1}}
\define@key{tablarara}{contenido}{\def\tablarara@contenido{#1}}
\define@key{tablarara}{caption}{\def\tablarara@caption{#1}}
\define@key{tablarara}{label}{\def\tablarara@label{#1}}

\presetkeys{tablarara}{
    titulo = {Título de la tabla},
    caption = {Descripción de la tabla},
    label = {anexo1:tab:tabla}
    contenido = {Contenido de la tabla},
}{}

\newcommand{\tablarara}[1]{
    \setkeys{tablarara}{#1}
    \begin{table}[H]
    \noindent \makebox[\textwidth] {
        \rowcolors{2}{Snow2}{Snow1}
        \begin{tabular}{ | l | p{10cm} | }
            \hline
            \rowcolor{Snow3}
            \tablarara@titulo \\
            \hline
            \tablarara@contenido
            \hline
        \end{tabular}
    }
    \caption{\tablarara@caption}
    \label{\tablarara@label}
    \end{table}
}

% Keys comunes para la mayoría de las tablas
\define@key{keygenerica}{id}{\def\keygenerica@id{#1}}
\define@key{keygenerica}{nombre}{\def\keygenerica@nombre{#1}}
\define@key{keygenerica}{fuentes}{\def\keygenerica@fuentes{#1}}
\define@key{keygenerica}{descripcion}{\def\keygenerica@descripcion{#1}}
\define@key{keygenerica}{importancia}{\def\keygenerica@importancia{#1}}
\define@key{keygenerica}{urgencia}{\def\keygenerica@urgencia{#1}}
\define@key{keygenerica}{estado}{\def\keygenerica@estado{#1}}
\define@key{keygenerica}{estabilidad}{\def\keygenerica@estabilidad{#1}}
\define@key{keygenerica}{comentarios}{\def\keygenerica@comentarios{#1}}

\presetkeys{keygenerica}{
    id = {0000},
    nombre = {},
    fuentes = {Propuesta inicial del trabajo de fin de grado},
    descripcion = {Descripción del elemento},
    importancia = {Vital},
    urgencia = {Inmediatamente},
    estado = {Validado},
    estabilidad = {Alta},
    comentarios = {No procede},
}{}

% Comando para definir un objetivo
\define@key{objetivo}{subobjetivos}{\def\objetivo@subobjetivos{#1}}

\presetkeys{objetivo}{
    subobjetivos = {No procede},
}{}

\newcommand{\objetivo}[2]{
    \setkeys{keygenerica}{#1}
    \setkeys{objetivo}{#2}
    \tablarara{
        titulo = {\textbf{OBJ-\keygenerica@id} & \textbf{\keygenerica@nombre}},
        contenido = {
            \textbf{Versión} & 1.0 (\textbf{TODO: FECHA}) \\
            \textbf{Autores} & Alberto Luengo Román \\
            \textbf{Fuentes} & \keygenerica@fuentes \\
            \textbf{Descripción} & \keygenerica@descripcion \\
            \textbf{Subobjetivos} & \objetivo@subobjetivos \\
            \textbf{Importancia} & \keygenerica@importancia \\
            \textbf{Urgencia} & \keygenerica@urgencia \\
            \textbf{Estado} & \keygenerica@estado \\
            \textbf{Estabilidad} & \keygenerica@estabilidad \\
            \textbf{Comentarios} & \keygenerica@comentarios \\
        },
        caption = {Objetivo \keygenerica@id: \keygenerica@nombre},
        label = {anexo1:tab:obj\keygenerica@id}
    }
}

% Comando para definir un requisito de información
\define@key{requisitoinfo}{objetivosAsociados}{\def\requisitoinfo@objetivosAsociados{#1}}
\define@key{requisitoinfo}{requisitosAsociados}{\def\requisitoinfo@requisitosAsociados{#1}}
\define@key{requisitoinfo}{datosEspecificos}{\def\requisitoinfo@datosEspecificos{#1}}
\define@key{requisitoinfo}{tiempoVidaMedio}{\def\requisitoinfo@tiempoVidaMedio{#1}}
\define@key{requisitoinfo}{tiempoVidaMaximo}{\def\requisitoinfo@tiempoVidaMaximo{#1}}
\define@key{requisitoinfo}{ocurrenciasSimultaneasMedio}{\def\requisitoinfo@ocurrenciasSimultaneasMedio{#1}}
\define@key{requisitoinfo}{ocurrenciasSimultaneasMaximo}{\def\requisitoinfo@ocurrenciasSimultaneasMaximo{#1}}

\presetkeys{requisitoinfo}{
    objetivosAsociados = {No procede},
    requisitosAsociados = {No procede},
    datosEspecificos = {No procede},
    tiempoVidaMedio = {No procede},
    tiempoVidaMaximo = {No procede},
    ocurrenciasSimultaneasMedio = {No procede},
    ocurrenciasSimultaneasMaximo = {No procede},
}{}

\newcommand{\requisitoinfo}[2]{
    \setkeys{keygenerica}{#1}
    \setkeys{requisitoinfo}{#2}

    \begin{table}[H]
    \noindent \makebox[\textwidth] {
        \rowcolors{2}{Snow2}{Snow1}
        \begin{tabular}{ | l | p{5cm} | p{5cm} | }
            \hline
            \rowcolor{Snow3}
            \textbf{RI-\keygenerica@id} & \multicolumn{2}{p{10.75cm}|}{\textbf{\keygenerica@nombre}} \\ 
            \hline
            \textbf{Versión} & \multicolumn{2}{p{10.75cm}|}{1.0 (\textbf{TODO: FECHA})} \\
            \textbf{Autores} & \multicolumn{2}{p{10.75cm}|}{Alberto Luengo Román} \\
            \textbf{Fuentes} & \multicolumn{2}{p{10.75cm}|}{\keygenerica@fuentes} \\
            \textbf{Objetivos asociados} & \multicolumn{2}{p{10.75cm}|}{\requisitoinfo@objetivosAsociados} \\
            \textbf{Requisitos asociados} & \multicolumn{2}{p{10.75cm}|}{\requisitoinfo@requisitosAsociados} \\
            \textbf{Descripción} & \multicolumn{2}{p{10.75cm}|}{\keygenerica@descripcion} \\
            \textbf{Datos específicos} & \multicolumn{2}{p{10.75cm}|}{\requisitoinfo@datosEspecificos} \\
                                                      & \textbf{Medio} & \textbf{Máximo} \\ \cline{2-3}
            \multirow{-2}{*}{\textbf{Tiempo de vida}} & \requisitoinfo@tiempoVidaMedio & \requisitoinfo@tiempoVidaMaximo \\
                                                               & \textbf{Medio} & \textbf{Máximo} \\ \cline{2-3}
            \multirow{-2}{*}{\textbf{Ocurrencias simultáneas}} & \requisitoinfo@ocurrenciasSimultaneasMedio & \requisitoinfo@ocurrenciasSimultaneasMaximo \\
            \textbf{Importancia} & \multicolumn{2}{p{10.75cm}|}{\keygenerica@importancia} \\
            \textbf{Urgencia} & \multicolumn{2}{p{10.75cm}|}{\keygenerica@urgencia} \\
            \textbf{Estado} & \multicolumn{2}{p{10.75cm}|}{\keygenerica@estado} \\
            \textbf{Estabilidad} & \multicolumn{2}{p{10.75cm}|}{\keygenerica@estabilidad} \\
            \textbf{Comentarios} & \multicolumn{2}{p{10.75cm}|}{\keygenerica@comentarios} \\
            \hline
        \end{tabular}
    }
    \caption{Requisito de información \keygenerica@id: \keygenerica@nombre}
    \label{anexo1:tab:ri\keygenerica@id}
    \end{table}
}

\makeatother
